\chapter{Laurent Series and the Residue Theorem}
\label{ch:b3:residues}

What happens to a \omterm{def:b3:holomorphic:holo}{holomorphic function} near a point where it
is \emph{not} defined? The answer is a \omterm{def:b3:complete:complete}{complete} trichotomy ---
removable point, pole, or essential singularity --- read off a
two-sided power series, the Laurent expansion. One coefficient
of that expansion, the \emph{\omterm{def:b3:residues:singularities}{residue}}, controls every \omterm{def:b3:holomorphic:contour}{contour
integral} around the singularity: the \omterm{def:b3:residues:singularities}{residue} theorem converts
hard definite integrals into finite algebra, counts zeros of
functions (argument principle, Rouch\'e), and proves the open
mapping theorem. We first upgrade Cauchy's theorem from
star-shaped domains to its definitive, homology-free form ---
Dixon's elegant argument --- so that all contours with
winding number zero around the complement become available.

\section{The global Cauchy theorem}

A \emph{cycle}\index{cycle} $\Gamma$ is a finite formal sum of
closed \omterm{def:b3:topology:pathconnected}{paths} $\gamma_1, \dots, \gamma_m$; integrals and
indices along $\Gamma$ are the corresponding sums, and
$\operatorname{im}\Gamma = \bigcup\operatorname{im}\gamma_j$.

\begin{theorem}[Cauchy, global form]
\label{thm:b3:residues:globalcauchy}
Let $\Omega \subseteq \C$ be open, $f \in \mathcal H(\Omega)$,
and $\Gamma$ a cycle in $\Omega$ such that
\[
\operatorname{Ind}_\Gamma(w) = 0
\qquad\text{for every } w \notin \Omega .
\]
Then, for all $z \in \Omega\setminus\operatorname{im}\Gamma$,
\[
\frac1{2\iu\pi}\int_\Gamma\frac{f(w)}{w - z}\,\dd w =
\operatorname{Ind}_\Gamma(z)\,f(z),
\qquad\text{and}\qquad
\int_\Gamma f(w)\,\dd w = 0 .
\]
\index{Cauchy's theorem (global)}
\end{theorem}

\begin{proof}[Proof (Dixon)]
Define $g \colon \Omega\times\Omega \to \C$ by
\[
g(z, w) = \begin{cases}
\dfrac{f(w) - f(z)}{w - z} & w \neq z,\\[4pt]
f'(z) & w = z .
\end{cases}
\]
\emph{$g$ is \omterm{def:b3:topology:continuity}{continuous}}: off the diagonal, clear. Near a
diagonal point $(a, a)$, expand $f$ in a power series at $a$
(\cref{thm:b3:holomorphic:analytic}): $f(w) - f(z) =
\sum_{n\geq1}c_n\bigl((w-a)^n - (z-a)^n\bigr)$, and dividing
each term by $w - z$ (factorization of $u^n - v^n$) gives, for
$z, w \in D(a, r)$,
\[
g(z, w) =
\sum_{n\geq1}c_n\sum_{j=0}^{n-1}(w-a)^{\,j}(z-a)^{\,n-1-j},
\]
valid also on the diagonal (each inner sum becomes
$n(z-a)^{n-1}$, summing to $f'(z)$). For $r$ small the series
converges uniformly on $D(a,r)^2$
($\abs{\text{term}} \leq n\abs{c_n}r^{n-1}$, summable inside
the radius): the sum is \omterm{def:b3:topology:continuity}{continuous}.

Set $h(z) = \frac1{2\iu\pi}\int_\Gamma g(z, w)\,\dd w$ on
$\Omega$: \omterm{def:b3:topology:continuity}{continuous} (uniform \omterm{def:b3:topology:continuity}{continuity} of $g$ on \omterm{def:b3:topology:compact}{compacts}),
and \omterm{def:b3:holomorphic:holo}{holomorphic} --- by Morera
(\cref{thm:b3:holomorphic:weierstrassconv}'s criterion): for a
triangle $T \subseteq \Omega$, Fubini gives $\int_{\partial
T}h = \frac1{2\iu\pi}\int_\Gamma\bigl(\int_{\partial T}g(z,
w)\dd z\bigr)\dd w = 0$, the inner integral vanishing because
$z \mapsto g(z, w)$ is \omterm{def:b3:holomorphic:holo}{holomorphic} on $\Omega$ (at $z = w$ the
singularity is removable: $g$ is \omterm{def:b3:topology:continuity}{continuous} there and
\omterm{def:b3:holomorphic:holo}{holomorphic} elsewhere --- the extension argument of
\cref{thm:b3:holomorphic:formula}'s proof).

On the \omterm{def:b3:topology:topology}{open set} $\Omega' = \{z \notin \operatorname{im}\Gamma
: \operatorname{Ind}_\Gamma(z) = 0\}$, define $h_1(z) =
\frac1{2\iu\pi}\int_\Gamma\frac{f(w)}{w - z}\dd w$:
\omterm{def:b3:holomorphic:holo}{holomorphic} on $\Omega'$ (Morera or differentiation under the
integral). For $z \in \Omega\cap\Omega'$:
\[
h(z) = \frac1{2\iu\pi}\int_\Gamma\frac{f(w)}{w-z}\dd w
- f(z)\operatorname{Ind}_\Gamma(z) = h_1(z) .
\]
By hypothesis $\Omega\cup\Omega' = \C$ ($w \notin \Omega
\Rightarrow \operatorname{Ind}_\Gamma(w) = 0$), so $h$ and
$h_1$ glue to an entire function $H$. Since the unbounded
component of the complement of $\operatorname{im}\Gamma$ lies
in $\Omega'$ and $h_1(z) \to 0$ as $\abs z \to \infty$ (ML
bound), $H$ is bounded and tends to $0$: Liouville
(\cref{cor:b3:holomorphic:liouville}) gives $H \equiv 0$. Thus
$h \equiv 0$ on $\Omega$, which is the integral formula.
Applying it, for fixed $a \in
\Omega\setminus\operatorname{im}\Gamma$, to $\tilde f(w) =
(w - a)f(w)$ at $z = a$:
\[
\frac1{2\iu\pi}\int_\Gamma f(w)\dd w =
\frac1{2\iu\pi}\int_\Gamma \frac{\tilde f(w)}{w - a}\dd w =
\operatorname{Ind}_\Gamma(a)\,\tilde f(a) = 0 .
\]
\end{proof}

\section{Laurent series and isolated singularities}

\begin{theorem}[Laurent expansion]\label{thm:b3:residues:laurent}
Let $f$ be \omterm{def:b3:holomorphic:holo}{holomorphic} on the annulus $A = \{r < \abs{z - a} <
R\}$ ($0 \leq r < R \leq \infty$). Then\index{Laurent series}
\[
f(z) = \sum_{n\in\Z}c_n\,(z - a)^n
\qquad\text{on } A,
\qquad
c_n = \frac1{2\iu\pi}\int_{C_\rho}\frac{f(w)}{(w -
a)^{n+1}}\,\dd w
\]
for any $r < \rho < R$ (independent of $\rho$), the two
half-series converging normally on \omterm{def:b3:topology:compact}{compact} subannuli. The
expansion is unique.
\end{theorem}

\begin{proof}
Fix $r < \rho_1 < \abs{z - a} < \rho_2 < R$ and let $\Gamma =
C_{\rho_2} - C_{\rho_1}$ (outer counterclockwise, inner
clockwise): a cycle in $A$ with
$\operatorname{Ind}_\Gamma(w) = 0$ for all $w \notin A$
(points inside the small disc: $1 - 1 = 0$; outside the big
one: $0 - 0$). By \cref{thm:b3:residues:globalcauchy},
$\operatorname{Ind}_\Gamma(z) = 1 - 0 = 1$ gives
\[
f(z) = \frac1{2\iu\pi}\int_{C_{\rho_2}}\frac{f(w)}{w - z}\dd
w - \frac1{2\iu\pi}\int_{C_{\rho_1}}\frac{f(w)}{w - z}\dd w .
\]
Expand the first kernel as in
\cref{thm:b3:holomorphic:analytic} (powers of $\frac{z -
a}{w - a}$, modulus $< 1$): the nonnegative part
$\sum_{n\geq0}c_n(z-a)^n$. In the second, expand the other
way: $\frac{-1}{w - z} = \frac1{(z-a)(1 - \frac{w - a}{z -
a})} = \sum_{m\geq0}\frac{(w-a)^m}{(z - a)^{m+1}}$, normally
convergent on $C_{\rho_1}$: the negative part
$\sum_{n\leq-1}c_n(z-a)^n$ with the stated coefficients (index
$n = -m-1$). Independence of $\rho$: the coefficient integrals
over $C_{\rho}$ and $C_{\rho'}$ differ by
$\int_\Gamma$ over a null-index cycle of the \omterm{def:b3:holomorphic:holo}{holomorphic}
$\frac{f(w)}{(w-a)^{n+1}}$ in $A$: zero, by
\cref{thm:b3:residues:globalcauchy} again. Uniqueness:
integrate $\sum c_n(z-a)^n$ against $(z - a)^{-m-1}$ over
$C_\rho$ term by term (normal convergence): only $n = m$
survives.
\end{proof}

\begin{definition}\label{def:b3:residues:singularities}
If $f$ is \omterm{def:b3:holomorphic:holo}{holomorphic} on a punctured disc $D(a, R)\setminus
\{a\}$, expand by Laurent ($r = 0$). Three exclusive
cases:\index{isolated singularity}
\begin{itemize}
  \item all $c_n = 0$ for $n < 0$: \emph{\omterm{thm:b3:residues:riemanncw}{removable
        singularity}} (the nonnegative series extends $f$
        \omterm{def:b3:holomorphic:holo}{holomorphically} to $a$);
  \item $c_n \neq 0$ for finitely many, at least one, $n < 0$:
        a \emph{pole} of order $m = -\min\{n : c_n \neq 0\}$;
        equivalently $f = g/(z-a)^m$, $g$ \omterm{def:b3:holomorphic:holo}{holomorphic}, $g(a)
        \neq 0$; equivalently $\abs{f(z)} \to \infty$ as $z\to
        a$;
  \item infinitely many negative $c_n \neq 0$:
        \emph{essential singularity}.
\end{itemize}
The \emph{residue}\index{residue} is
$\operatorname{Res}(f, a) = c_{-1}$. A function \omterm{def:b3:holomorphic:holo}{holomorphic} on
$\Omega$ minus a set of poles is
\emph{meromorphic}\index{meromorphic function} on $\Omega$.
\end{definition}

\begin{theorem}[Riemann; Casorati--Weierstrass]
\label{thm:b3:residues:riemanncw}
Let $f$ be \omterm{def:b3:holomorphic:holo}{holomorphic} on $D(a,R)\setminus\{a\}$.
\begin{enumerate}
  \item (Riemann) If $f$ is \emph{bounded} near $a$, the
        singularity is removable.\index{removable singularity}
  \item (Casorati--Weierstrass) If $a$ is essential, then
        $f\bigl(D(a,\varepsilon)\setminus\{a\}\bigr)$ is dense
        in $\C$ for every $\varepsilon$.
        \index{Casorati--Weierstrass theorem}
\end{enumerate}
\end{theorem}

\begin{proof}
(1) For $n < 0$ and $\rho \to 0$: $\abs{c_n} \leq
M\rho^{-n-1}\cdot\rho\cdot\rho^{-\,n}\dots$ by ML on
$C_\rho$: $\abs{c_n} \leq \frac{1}{2\pi}\,2\pi\rho\cdot
M\rho^{-(n+1)} = M\rho^{-n} \to 0$ (as $-n > 0$): all
negative coefficients vanish.
(2) If some value $b$ were not approached: $\abs{f - b} \geq
\delta$ near $a$, so $g = 1/(f - b)$ is \omterm{def:b3:holomorphic:holo}{holomorphic} and
bounded near $a$: removable (1), $g$ extends with value $c$.
If $c \neq 0$, $f = b + 1/g$ is bounded near $a$: removable
--- excluded. If $c = 0$, $g$ has a zero of finite order $m$
at $a$ (\cref{thm:b3:holomorphic:identity}; $g \not\equiv 0$),
and $f = b + 1/g$ has a pole of order $m$: excluded again.
\end{proof}

\section{The residue theorem}

\begin{theorem}[Residue theorem]\label{thm:b3:residues:residue}
Let $\Omega$ be open, $S \subseteq \Omega$ finite, $f \in
\mathcal H(\Omega\setminus S)$, and $\Gamma$ a cycle in
$\Omega\setminus S$ with $\operatorname{Ind}_\Gamma(w) = 0$
for every $w \notin \Omega$. Then\index{residue theorem}
\[
\frac{1}{2\iu\pi}\int_\Gamma f(z)\,\dd z
= \sum_{a\in S}\operatorname{Ind}_\Gamma(a)\,
\operatorname{Res}(f, a) .
\]
\end{theorem}

\begin{proof}
For each $a \in S$, let $P_a(z) = \sum_{n\leq-1}c_n^{(a)}(z -
a)^n$ be the principal part of $f$ at $a$: a series
converging on $\C\setminus\{a\}$ (its radius in $1/(z-a)$ is
infinite: the Laurent tail converges for all small $\abs{z-a}$
hence, being a power series in $(z-a)^{-1}$, everywhere), and
\omterm{def:b3:holomorphic:holo}{holomorphic} there. Then $g = f - \sum_{a\in S}P_a$ has
removable singularities at each point of $S$ (its Laurent
expansion at $a$ has no negative part: the other $P_{a'}$ are
\omterm{def:b3:holomorphic:holo}{holomorphic} at $a$), so $g$ extends \omterm{def:b3:holomorphic:holo}{holomorphically} to
$\Omega$, and \cref{thm:b3:residues:globalcauchy} gives
$\int_\Gamma g = 0$. It remains to integrate each $P_a$:
term-by-term (normal convergence on the \omterm{def:b3:topology:compact}{compact}
$\operatorname{im}\Gamma$, which avoids $a$),
\[
\frac1{2\iu\pi}\int_\Gamma(z - a)^n\,\dd z = 0 \ (n \leq -2:
\text{primitive } \tfrac{(z-a)^{n+1}}{n+1}),
\qquad
\frac1{2\iu\pi}\int_\Gamma\frac{\dd z}{z - a} =
\operatorname{Ind}_\Gamma(a),
\]
so $\frac1{2\iu\pi}\int_\Gamma P_a =
c_{-1}^{(a)}\operatorname{Ind}_\Gamma(a)$. Sum over $a$.
\end{proof}

\begin{method}[Computing residues]\label{met:b3:residues:compute}
Simple pole: $\operatorname{Res}(f, a) = \lim_{z\to a}(z -
a)f(z)$; for $f = g/h$ with $g(a) \neq 0$, $h(a) = 0$, $h'(a)
\neq 0$: $\operatorname{Res} = g(a)/h'(a)$. Pole of order $m$:
$\operatorname{Res}(f, a) = \frac1{(m-1)!}\lim_{z\to
a}\bigl((z-a)^mf(z)\bigr)^{(m-1)}$. Essential singularities:
expand and read $c_{-1}$ (e.g.\ from known series). Always
verify which poles the contour actually encircles, and with
which index.
\end{method}

\begin{example}[The four classical integral types]
\label{ex:b3:residues:integrals}
(a) \emph{Rational over $\R$}: for $\int_\R\frac{\dd x}{1 +
x^4}$, close with a large semicircle $S_R$ in the upper
half-plane: the integrand is $O(R^{-4})$ there, so
$\int_{S_R} \to 0$ (ML), and the \omterm{def:b3:residues:singularities}{residue} theorem with the
poles $\eu^{\iu\pi/4}, \eu^{3\iu\pi/4}$ (simple, \omterm{def:b3:residues:singularities}{residues}
$\frac1{4z^3} = \frac{z}{4z^4} = -\frac z4$ at a pole) gives
\[
\int_\R\frac{\dd x}{1 + x^4}
= 2\iu\pi\Bigl(-\frac{\eu^{\iu\pi/4}}4 -
\frac{\eu^{3\iu\pi/4}}4\Bigr)
= \frac{\pi}{\sqrt2} .
\]
(b) \emph{Fourier type}: for $t \geq 0$,
$\int_\R\frac{\eu^{\iu tx}}{1 + x^2}\dd x =
2\iu\pi\operatorname{Res}\bigl(\tfrac{\eu^{\iu tz}}{1+z^2},
\iu\bigr) = 2\iu\pi\frac{\eu^{-t}}{2\iu} = \pi\eu^{-t}$ ---
the upper semicircle works because $\abs{\eu^{\iu tz}} =
\eu^{-t\operatorname{Im}z} \leq 1$ there; taking real parts:
$\int_\R\frac{\cos(tx)}{1+x^2}\dd x = \pi\eu^{-\abs t}$,
settling \cref{exo:b3:lebesgue:10}'s admitted formula.
(c) \emph{Trigonometric over a period}: substitute $z =
\eu^{\iu t}$, $\cos t = \frac{z + z^{-1}}2$, $\dd t =
\frac{\dd z}{\iu z}$: $\int_0^{2\pi}\frac{\dd t}{a + \cos t}$
($a > 1$) becomes a \omterm{def:b3:residues:singularities}{residue} count inside the unit circle
(\cref{exo:b3:residues:2}).
(d) \emph{Series}: pair $f$ with $\pi\cot(\pi z)$, whose poles
are the integers with \omterm{def:b3:residues:singularities}{residue} $1$: the weekend problem sums
$\sum n^{-2}$ and $\sum n^{-4}$ this way.
\end{example}

\begin{omfigure}
\begin{tikzpicture}[scale=1.15]
  \draw[->] (-2.5,0) -- (2.6,0) node[right, font=\small]
    {$\operatorname{Re}$};
  \draw[->] (0,-0.5) -- (0,2.5) node[above, font=\small]
    {$\operatorname{Im}$};
  \draw[omThm, very thick, ->] (-2,0) -- (0.3,0);
  \draw[omThm, very thick] (0.3,0) -- (2,0);
  \draw[omThm, very thick] (2,0) arc (0:90:2);
  \draw[omThm, very thick, ->] (2,0) arc (0:50:2);
  \draw[omThm, very thick] (0,2) arc (90:180:2);
  \node[omThm, font=\small] at (1.75,1.55) {$S_R$};
  \node[font=\small] at (2.2,-0.25) {$R$};
  \node[font=\small] at (-2.15,-0.25) {$-R$};
  \fill[omDef] (45:1.19) circle (1.6pt)
    node[above right, font=\small] {$\eu^{\iu\pi/4}$};
  \fill[omDef] (135:1.19) circle (1.6pt)
    node[above left, font=\small] {$\eu^{3\iu\pi/4}$};
  \fill[black!40] (-45:1.19) circle (1.6pt);
  \fill[black!40] (-135:1.19) circle (1.6pt);
\end{tikzpicture}

{\small The semicircle contour for $\int_\R\frac{\dd x}{1 +
x^4}$: as $R \to \infty$ the arc contributes $O(R^{-3})$, and
the \omterm{def:b3:residues:singularities}{residue} theorem counts the two enclosed poles (blue). The
two lower poles (gray) are outside: index $0$.}
\end{omfigure}

\section{The argument principle and Rouch\'e's theorem}

\begin{theorem}[Argument principle]
\label{thm:b3:residues:argument}
Let $f$ be \omterm{def:b3:residues:singularities}{meromorphic} on $\Omega$, with zeros $z_j$ (orders
$m_j$) and poles $p_k$ (orders $\mu_k$), and $\gamma$ a closed
\omterm{def:b3:topology:pathconnected}{path} in $\Omega$ avoiding them all, with
$\operatorname{Ind}_\gamma = 0$ off $\Omega$. Then
\index{argument principle}
\[
\frac1{2\iu\pi}\int_\gamma\frac{f'(z)}{f(z)}\,\dd z
= \sum_j m_j\operatorname{Ind}_\gamma(z_j)
- \sum_k \mu_k\operatorname{Ind}_\gamma(p_k)
\]
(finitely many terms are nonzero). For a simple
counterclockwise contour, the integral counts zeros minus
poles inside, with multiplicity --- and equals the winding
number of the image \omterm{def:b3:topology:pathconnected}{path} $f\circ\gamma$ around $0$.
\end{theorem}

\begin{proof}
Near a zero of order $m$: $f = (z-a)^mg$, $g(a) \neq 0$, so
$\frac{f'}f = \frac m{z - a} + \frac{g'}g$ with the second
term \omterm{def:b3:holomorphic:holo}{holomorphic} near $a$: a simple pole of \omterm{def:b3:residues:singularities}{residue} $m$. Near
a pole of order $\mu$: $f = (z-a)^{-\mu}g$ gives \omterm{def:b3:residues:singularities}{residue}
$-\mu$. Elsewhere $\frac{f'}f$ is \omterm{def:b3:holomorphic:holo}{holomorphic}. (The zeros and
poles with nonzero index lie in a \omterm{def:b3:topology:compact}{compact} region encircled by
$\gamma$; by the identity theorem they are finite in number
there, $f \not\equiv 0$.) Apply
\cref{thm:b3:residues:residue}. The last remark:
$\frac1{2\iu\pi}\int_\gamma\frac{f'}f =
\frac1{2\iu\pi}\int_{f\circ\gamma}\frac{\dd w}w =
\operatorname{Ind}_{f\circ\gamma}(0)$ (substitute $w =
f(\gamma(t))$).
\end{proof}

\begin{theorem}[Rouch\'e]\label{thm:b3:residues:rouche}
Let $f, g$ be \omterm{def:b3:holomorphic:holo}{holomorphic} on $\Omega$, and $\gamma$ a closed
\omterm{def:b3:topology:pathconnected}{path} with $\operatorname{Ind}_\gamma \in \{0,1\}$, zero off
$\Omega$ (a simple contour). If\index{Rouch\'e's theorem}
\[
\abs{g(z)} < \abs{f(z)} \qquad \text{on }
\operatorname{im}\gamma,
\]
then $f$ and $f + g$ have the same number of zeros (with
multiplicity) in the region $\{\operatorname{Ind}_\gamma =
1\}$.
\end{theorem}

\begin{proof}
For $t \in \intcc01$, $f_t = f + tg$ has no zero on
$\operatorname{im}\gamma$ ($\abs{f_t} \geq \abs f - \abs g >
0$), so
\[
N(t) = \frac1{2\iu\pi}\int_\gamma
\frac{f_t'(z)}{f_t(z)}\,\dd z
\]
is well defined; it counts the zeros in the enclosed region
(\cref{thm:b3:residues:argument}; no poles). $N$ is
\omterm{def:b3:topology:continuity}{continuous} in $t$ (the integrand is jointly \omterm{def:b3:topology:continuity}{continuous}, the
denominators uniformly bounded below --- dominated
convergence) and integer-valued: constant. $N(0) = N(1)$.
\end{proof}

\begin{corollary}[Open mapping theorem]
\label{cor:b3:residues:openmapping}
A nonconstant \omterm{def:b3:holomorphic:holo}{holomorphic function} on a \omterm{def:b3:topology:connected}{connected} \omterm{def:b3:topology:topology}{open set} is
an open map. In particular (again) the maximum principle
holds, and a \omterm{def:b3:holomorphic:holo}{holomorphic} bijection has \omterm{def:b3:holomorphic:holo}{holomorphic}
inverse.\index{open mapping theorem (holomorphic)}
\end{corollary}

\begin{proof}
Let $f(a) = b$; $f - b$ has a zero of some finite order $m
\geq 1$ at $a$ (identity theorem: $f \not\equiv b$). Choose
$r$ with $f - b$ zero-free on $\bar D(a, r)\setminus\{a\}$
(\omterm{thm:b3:holomorphic:identity}{isolated zeros}) and let $\delta = \min_{\abs{z - a} =
r}\abs{f(z) - b} > 0$. For $\abs{w - b} < \delta$: on the
circle, $\abs{(b - w)} < \delta \leq \abs{f - b}$, so
Rouch\'e ($f - b$ versus constant $b - w$) says $f - w$ has
exactly $m$ zeros in $D(a, r)$: every such $w$ is attained
--- $f(D(a,r)) \supseteq D(b, \delta)$: open. Maximum
principle: an \omterm{def:b3:topology:interior}{interior} maximum of $\abs f$ is impossible for
a nonconstant $f$, its image around $f(a)$ containing points
of larger modulus. Inverse: a \omterm{def:b3:holomorphic:holo}{holomorphic} bijection $f$ is
open, so $f^{-1}$ is \omterm{def:b3:topology:continuity}{continuous}; the zero of $f - f(a)$ at
$a$ is simple ($m \geq 2$ would give $m$ preimages of nearby
values --- distinct ones, since $f'$ vanishes only at
isolated points, so near $a$ the $m$ zeros of $f - w$ are
simple and distinct for generic small $w$: contradiction
with injectivity); then $f'(a) \neq 0$ and the difference
quotient of $f^{-1}$ converges: $\bigl(f^{-1}\bigr)'(b) =
1/f'(a)$.
\end{proof}

\section{Exercises}

\begin{exercise}[$\star$]\label{exo:b3:residues:1}
Classify the singularity at $0$ and compute the \omterm{def:b3:residues:singularities}{residue}:
\[
\frac{\sin z}{z},\qquad
\frac{\eu^z - 1}{z^2},\qquad
\frac{1}{z(z-1)^2},\qquad
\frac{\cos z}{z^3},\qquad
\eu^{1/z},\qquad
\frac1{\sin z} .
\]
Also give the \omterm{def:b3:residues:singularities}{residue} of the third at $z = 1$ and of the last
at $z = \pi$.
\end{exercise}

\begin{exercise}[$\star$]\label{exo:b3:residues:2}
For $a > 1$ compute, via $z = \eu^{\iu t}$:
\[
\int_0^{2\pi}\frac{\dd t}{a + \cos t}
= \frac{2\pi}{\sqrt{a^2 - 1}} .
\]
Check the limit behaviors $a \to 1^+$ and $a \to \infty$.
\end{exercise}

\begin{exercise}[$\star\star$]\label{exo:b3:residues:3}
Compute with semicircle contours, justifying the arc estimates:
\[
\int_\R\frac{x^2}{1 + x^6}\,\dd x = \frac\pi3,
\qquad
\int_\R\frac{\dd x}{(1 + x^2)^{2}} = \frac\pi2
\quad\text{(re-deriving \cref{exo:b3:fouriertransform:3})}.
\]
\end{exercise}

\begin{exercise}[$\star\star$]\label{exo:b3:residues:4}
Prove, for $t \geq 0$ and $a > 0$:
\[
\int_\R\frac{\cos(tx)}{x^2 + a^2}\,\dd x =
\frac{\pi}{a}\,\eu^{-at},
\]
and deduce the Fourier transform of $x \mapsto \eu^{-a\abs
x}$ by inversion --- comparing with
\cref{exo:b3:fouriertransform:1}.
\end{exercise}

\begin{exercise}[$\star\star$]\label{exo:b3:residues:5}
Expand $f(z) = \dfrac1{(z-1)(z-2)}$ in \omterm{thm:b3:residues:laurent}{Laurent series} in each
of the three regions $\abs z < 1$, $1 < \abs z < 2$, $\abs z >
2$. Why do the three expansions differ? Explain why the
coefficient of $z^{-1}$ in the second and third expansions is
\emph{not} a \omterm{def:b3:residues:singularities}{residue} of $f$ at $0$ ($f$ has no singularity
there), and compute the actual \omterm{def:b3:residues:singularities}{residues} of $f$, at $1$ and at
$2$.
\end{exercise}

\begin{exercise}[$\star\star$]\label{exo:b3:residues:6}
(a) Show that $\eu^{1/z}$ has an essential singularity at $0$
and verify Casorati--Weierstrass by hand: solve $\eu^{1/z} =
w$ explicitly for any $w \neq 0$, exhibiting solutions
arbitrarily close to $0$.
(b) Show that $\abs{\eu^{1/z}}$ is unbounded on every
punctured \omterm{def:b3:topology:topology}{neighborhood} of $0$ yet $\eu^{1/z}$ has no pole:
which limit fails?
\end{exercise}

\begin{exercise}[$\star\star$]\label{exo:b3:residues:7}
Count with Rouch\'e:
(a) the zeros of $z^7 - 4z^3 + z - 1$ in $\abs z < 1$;
(b) the zeros of $z^4 + 5z + 1$ in $\abs z < 1$ and in $1 <
\abs z < 2$;
(c) re-prove d'Alembert--Gauss: a monic degree-$n$ polynomial
has $n$ zeros in some large disc \emph{(compare with
$z^n$)}.
\end{exercise}

\begin{exercise}[$\star\star\star$]\label{exo:b3:residues:8}
(Hurwitz) Let $f_n \to f$ uniformly on \omterm{def:b3:topology:compact}{compacts}, $f_n \in
\mathcal H(\Omega)$, $\Omega$ \omterm{def:b3:topology:connected}{connected}, $f \not\equiv 0$.
(a) Show that if all $f_n$ are zero-free, so is $f$.
\emph{(If $f(a) = 0$: argument principle on a small circle
around $a$, and \cref{thm:b3:holomorphic:weierstrassconv} to
pass to the limit in $\int f_n'/f_n$.)}
(b) Show that if all $f_n$ are injective, $f$ is injective or
constant. \emph{(Apply (a) to $z \mapsto f_n(z) - f_n(w)$ on
$\Omega\setminus\{w\}$.)}
\end{exercise}

\begin{exercise}[$\star\star\star$]\label{exo:b3:residues:9}
For $n \geq 2$, integrate $\frac1{1 + z^n}$ over the boundary
of the sector $\{0 \leq \arg z \leq \frac{2\pi}n,\ \abs z
\leq R\}$ and deduce
\[
\int_0^{+\infty}\frac{\dd x}{1 + x^n} =
\frac{\pi}{n\,\sin(\pi/n)} .
\]
Verify $n = 2$ against $\arctan$, and the limit $n \to
\infty$.
\end{exercise}

\begin{exercise}[$\star\star$]\label{exo:b3:residues:10}
Let $f$ be a rational function with $\deg(\text{denominator})
\geq \deg(\text{numerator}) + 2$. Show that the sum of
\emph{all} \omterm{def:b3:residues:singularities}{residues} of $f$ is zero \emph{(integrate over
larger and larger circles)}. Use this to recompute the partial
fraction decomposition of $\frac1{z(z-1)(z-2)}$ with no linear
algebra.
\end{exercise}

\begin{exercise}[$\star\star\star$]\label{exo:b3:residues:11}
(The keyhole: Euler's reflection integral) For $0 < a < 1$,
compute
\[
I(a) = \int_0^{\infty}\frac{x^{a-1}}{1 + x}\,\dd x =
\frac{\pi}{\sin(\pi a)}
\]
by integrating $f(z) = \frac{z^{a-1}}{1+z} =
\frac{\eu^{(a-1)\log z}}{1 + z}$ (logarithm cut along
$\R_+$, $\arg z \in \intoo0{2\pi}$) over the keyhole
contour: out along the top of the cut from $\varepsilon$ to
$R$, around $C_R$, back under the cut, around
$C_\varepsilon$. Justify: the two straight stretches differ
by the factor $\eu^{2\iu\pi(a-1)}$, the circle contributions
vanish ($R^{a-1}\cdot R \to 0$ and
$\varepsilon^{a-1}\cdot\varepsilon \to 0$), and the unique
pole $z = -1$ has \omterm{def:b3:residues:singularities}{residue} $\eu^{\iu\pi(a - 1)}$. Deduce also
$\Gamma(a)\Gamma(1 - a) = \frac\pi{\sin\pi a}$ \emph{(write
$\Gamma(a)\Gamma(1-a) = B(a, 1-a)$ by
\cref{pb:b3:lebesgue:1} and substitute $t =
\frac{x}{1+x}$)}.
\end{exercise}

\begin{exercise}[$\star\star$]\label{exo:b3:residues:12}
(Counting zeros with the argument principle, numerically)
Let $P(z) = z^4 + 8z + 1$.
(a) How many zeros in the unit disc? \emph{(Rouch\'e against
$8z + 1$.)}
(b) How many in the annulus $1 < \abs z < 3$?
\emph{(Rouch\'e against $z^4$ on $\abs z = 3$.)} Sharpen:
show every zero has modulus $< 2.1$.
(c) How many in the right half-plane? \emph{(Count on
$\abs z = 2$ first; then track the image of the imaginary
axis: $P(\iu t) = t^4 + 1 + 8\iu t$ has positive real part
throughout, so no zeros on the axis, and the argument
variation along it is computable --- conclude with a large
half-disc.)}
\end{exercise}

\section{Problem: \texorpdfstring{$\zeta(2k)$}{zeta(2k)} by the
cotangent}

\begin{problem}[{Weekend problem --- summing $\sum n^{-2k}$
with \omterm{def:b3:residues:singularities}{residues}}]\label{pb:b3:residues:1}
The \omterm{def:b3:residues:singularities}{residue} theorem sums series: pairing a rational function
with $\pi\cot(\pi z)$, whose poles sit at the integers, turns
$\sum_{n}f(n)$ into a \omterm{def:b3:residues:singularities}{residue} count. We prove the method and
compute $\zeta(2) = \frac{\pi^2}{6}$ and $\zeta(4) =
\frac{\pi^4}{90}$ --- the values found by Fourier series in
Year 2 and by operator traces in \cref{ch:b3:spectral}, now by
contour integration.

\medskip
\noindent\textbf{Part I --- The cotangent kernel.}
\begin{enumerate}
  \item Show that $\pi\cot(\pi z)$ is \omterm{def:b3:residues:singularities}{meromorphic} on $\C$ with
        simple poles exactly at $z = n \in \Z$, each of
        \omterm{def:b3:residues:singularities}{residue} $1$ \emph{(compute
        $\lim_{z\to n}(z-n)\pi\cot\pi z$)}.
  \item Compute the beginning of the Laurent expansion at $0$:
        \[
        \pi\cot(\pi z) = \frac1z - \frac{\pi^2}{3}\,z -
        \frac{\pi^4}{45}\,z^3 + O(z^5) ,
        \]
        by dividing the power series of $\cos$ by that of
        $\sin$ (justify the division: $\frac{\sin\pi z}{\pi z}$
        is \omterm{def:b3:holomorphic:holo}{holomorphic} and nonzero near $0$, so its
        reciprocal is \omterm{def:b3:holomorphic:holo}{holomorphic}; identify coefficients
        through order $3$).
  \item Let $C_N$ be the boundary of the square with vertices
        $(\pm1\pm\iu)(N + \frac12)$. Show that
        $\abs{\cot(\pi z)} \leq 2$ on $C_N$ for every $N \geq
        1$. \emph{(On vertical sides, $\cot(\pi(\pm(N +
        \frac12) + \iu y)) = \mp\tan(\iu\pi y)$, of modulus
        $\abs{\tanh(\pi y)} \leq 1$; on horizontal sides
        $\abs y = N + \frac12$, bound
        $\abs{\cot(\pi(x\pm\iu y))} \leq \coth(\pi y) \leq
        \coth(\pi/2) < 1.1$.)}
\end{enumerate}

\medskip
\noindent\textbf{Part II --- The summation theorem.}
\begin{enumerate}[resume]
  \item Let $f$ be rational, \omterm{def:b3:holomorphic:holo}{holomorphic} at the integers,
        with $\deg(\text{denom}) \geq \deg(\text{num}) + 2$.
        Using the \omterm{def:b3:residues:singularities}{residue} theorem on $C_N$ and the bound of
        question 3, prove:
        \[
        \lim_{N\to\infty}\ \sum_{n = -N}^{N} f(n)
        = -\sum_{p\ \text{pole of}\ f}
        \operatorname{Res}\bigl(\pi\cot(\pi z)f(z),\,p\bigr).
        \]
  \item Where does the argument need the degree condition?
        Show by example (take $f(z) = 1/(z + \frac12)$)
        that for slower decay the symmetric limit
        may still exist while the two-sided series diverges
        --- and that the formula then computes the
        \emph{principal value}.
\end{enumerate}

\medskip
\noindent\textbf{Part III --- The values.}
\begin{enumerate}[resume]
  \item Apply the method to $f(z) = 1/z^2$: here $f$ has its
        pole \emph{at} an integer, so run the argument
        directly --- integrate $g(z) = \frac{\pi\cot(\pi
        z)}{z^2}$ over $C_N$, show the integral $\to 0$, and
        compute $\operatorname{Res}(g, 0)$ from question 2.
        Conclude:
        \[
        2\sum_{n\geq1}\frac1{n^2} = \frac{\pi^2}{3},
        \qquad \zeta(2) = \frac{\pi^2}6 .
        \]
  \item Same with $g(z) = \frac{\pi\cot(\pi z)}{z^4}$:
        compute $\operatorname{Res}(g, 0)$ and deduce
        $\zeta(4) = \frac{\pi^4}{90}$.
  \item Explain the general pattern: for every $k \geq 1$,
        $\zeta(2k)$ is $-\frac12$ times the coefficient of
        $z^{2k-1}$ in the Laurent expansion of $\pi\cot(\pi
        z)$ at $0$ --- a rational multiple of $\pi^{2k}$.
        Compute $\zeta(6)$ by pushing question 2's division
        one step further. What does the method say about
        $\zeta(3)$ --- and why does it say nothing?
\end{enumerate}

\medskip
\noindent\textbf{Part IV --- The partial fraction expansion
of the cotangent.}
\begin{enumerate}[resume]
  \item Fix $w \in \C\setminus\Z$ and apply the method of
        Part II to $f(z) = \dfrac{1}{(z - w)(z + w)}$ ---
        noting that $\pi\cot(\pi z)f(z)$ now has additional
        simple poles at $\pm w$, whose \omterm{def:b3:residues:singularities}{residues} must join the
        count. Deduce the partial fraction expansion
        \[
        \pi\cot(\pi w) = \frac1w +
        \sum_{n\geq1}\frac{2w}{w^2 - n^2},
        \]
        the series converging normally on \omterm{def:b3:topology:compact}{compact} subsets of
        $\C\setminus\Z$.
  \item Recover from this expansion, by expanding each term
        in powers of $w$ (justify the interchange), the
        \emph{same} Laurent coefficients as in question 2 ---
        the circle closes: Euler's formula
        $\sum\frac1{n^2} = \frac{\pi^2}6$ is the coefficient
        of $w$ in the cotangent's two faces. Compare with the
        Fourier-series proof (Year 2) and the trace proof
        (\cref{pb:b3:spectral:1}): three theories, one
        number.
\end{enumerate}

\medskip
\noindent\textbf{Part V --- Euler's product for the sine.}
The expansion of question 9 is the logarithmic derivative of
an infinite product; we now prove Euler's 1734 factorization
honestly.
\begin{enumerate}[resume]
  \item For $N \geq 1$ set $P_N(z) =
        z\prod_{n=1}^{N}\bigl(1 - \frac{z^2}{n^2}\bigr)$.
        Show that $P_N$ converges, uniformly on every disc
        $\bar D(0, R)$, to an entire function $P$ whose zeros
        are exactly the integers, all simple. \emph{(For $n
        \geq 2R$ write the factor as $\exp\log(1 -
        z^2/n^2)$ with the principal logarithm of
        \cref{exo:b3:holomorphic:3}, bound $\abs{\log(1+u)}
        \leq 2\abs u$ for $\abs u \leq \frac12$ via the
        series, and exponentiate the normally convergent sum
        of logarithms; the finitely many remaining factors
        are a polynomial. Conclude with
        \cref{thm:b3:holomorphic:weierstrassconv}.)}
  \item Show that on $\C\setminus\Z$,
        \[
        \frac{P'(z)}{P(z)} = \frac1z +
        \sum_{n\geq1}\frac{2z}{z^2 - n^2} = \pi\cot(\pi z)
        \]
        \emph{(differentiate the finite products, pass to
        the limit using
        \cref{thm:b3:holomorphic:weierstrassconv} and the
        zero-freeness of $P$ off $\Z$, and quote question
        9)}.
  \item Show that $Q = \sin(\pi z)/P(z)$ extends to a
        zero-free entire function with $Q' = 0$, and
        conclude \emph{Euler's product}:
        \[
        \sin(\pi z) = \pi z\prod_{n\geq1}
        \Bigl(1 - \frac{z^2}{n^2}\Bigr)
        \qquad (z \in \C) .
        \]
  \item (Wallis, 1655) Evaluate at $z = \frac12$:
        \[
        \frac\pi2 = \prod_{n\geq1}\frac{4n^2}{4n^2 - 1}
        = \lim_{N\to\infty}
        \frac{2\cdot2\cdot4\cdot4\cdots(2N)(2N)}
        {1\cdot3\cdot3\cdot5\cdots(2N-1)(2N+1)} .
        \]
  \item For $\abs z < 1$, expand the logarithm of the
        product as a double series (justify the
        rearrangement) and recover $\zeta(2) =
        \frac{\pi^2}6$ by matching the coefficient of $z^3$
        in $\sin(\pi z) = \pi z - \frac{\pi^3}6z^3 + \cdots$
        --- the product's face of Euler's number.
\end{enumerate}

\medskip
\noindent\textbf{Part VI --- Sister kernels.} The cotangent
has siblings; each prices its own family of series.
\begin{enumerate}[resume]
  \item Differentiate question 9's expansion term by term
        (justified by
        \cref{thm:b3:holomorphic:weierstrassconv}) to
        obtain, normally on \omterm{def:b3:topology:compact}{compacts} of $\C\setminus\Z$,
        \[
        \frac{\pi^2}{\sin^2(\pi z)} =
        \sum_{n\in\Z}\frac1{(z - n)^2} .
        \]
  \item Evaluate at $z = \frac12$:
        $\sum_{m\geq0}\frac1{(2m+1)^2} = \frac{\pi^2}8$;
        recover $\zeta(2)$ once more by splitting the
        integers by parity.
  \item Verify the duplication identity $\tan\theta =
        \cot\theta - 2\cot(2\theta)$ and deduce
        \[
        \pi\tan(\pi z) =
        \sum_{m\geq0}\frac{8z}{(2m+1)^2 - 4z^2} ,
        \]
        normally on \omterm{def:b3:topology:compact}{compacts} avoiding $\frac12 + \Z$.
  \item Expand around $0$ ($\abs z < \frac12$; Fubini
        again): with $\lambda(s) =
        \sum_{m\geq0}(2m+1)^{-s}$,
        \[
        \pi\tan(\pi z) =
        \sum_{k\geq0}8\cdot4^k\,\lambda(2k+2)\,z^{2k+1} ;
        \]
        compare with $\tan u = u + \frac{u^3}3 + O(u^5)$ to
        recover $\lambda(2) = \frac{\pi^2}8$ and to get
        $\lambda(4) = \frac{\pi^4}{96}$, then cross-check
        $\zeta(4) = \frac{\pi^4}{90}$ via $\lambda(4) = (1 -
        2^{-4})\,\zeta(4)$.
  \item Verify $\frac1{\sin\theta} = \cot\frac\theta2 -
        \cot\theta$ and deduce
        \[
        \frac{\pi}{\sin(\pi z)} = \frac1z +
        \sum_{n\geq1}(-1)^n\,\frac{2z}{z^2 - n^2} .
        \]
        Check the signs against the \omterm{def:b3:residues:singularities}{residues} of
        $\pi/\sin(\pi z)$ at the integers.
  \item Read off the coefficient of $z$: $\eta(2) =
        \sum_{n\geq1}\frac{(-1)^{n-1}}{n^2} =
        \frac{\pi^2}{12}$, and confirm the consistency
        $\eta(2) = (1 - 2^{1-2})\,\zeta(2)$.
  \item (Finale) Evaluate question 20's expansion at $z =
        \frac12$ and deduce \emph{Leibniz's formula}
        \[
        \frac\pi4 = 1 - \frac13 + \frac15 - \frac17 +
        \cdots
        \]
        Close with a short paragraph: one kernel per
        arithmetic --- which kernel prices which family of
        series, and why all of them are structurally blind
        to $\zeta(3)$.

\end{enumerate}
\medskip
\noindent\textbf{Part VII --- The full price list: Bernoulli
numbers.}
\begin{enumerate}[resume]
  \item Combine the partial fraction expansion of $\pi
        z\cot(\pi z)$ with the generating function of the
        Bernoulli numbers ($\frac{w}{\eu^w - 1} =
        \sum_n\frac{B_n}{n!}w^n$,
        \cref{pb:b3:holomorphic:1}, Part VI): from
        \[
        \pi z\cot(\pi z) = \iu\pi z +
        \frac{2\iu\pi z}{\eu^{2\iu\pi z} - 1}
        \]
        (prove this identity first), deduce the closed form
        \[
        \zeta(2k) = (-1)^{k+1}\,
        \frac{(2\pi)^{2k}\,B_{2k}}{2\,(2k)!}
        \qquad (k \geq 1).
        \]
  \item Verify the formula against $B_2 = \frac16$, $B_4 =
        -\frac1{30}$, $B_6 = \frac1{42}$: recover $\zeta(2) =
        \frac{\pi^2}6$, $\zeta(4) = \frac{\pi^4}{90}$, and
        compute $\zeta(6) = \frac{\pi^6}{945}$.
  \item (Euler's recursion) Expand both sides of
        $\bigl(z\cot z\bigr)' = \cot z - z(1 + \cot^2z)$ ---
        or square the cotangent series directly --- to prove
        \[
        \Bigl(k + \frac12\Bigr)\zeta(2k) =
        \sum_{j=1}^{k-1}\zeta(2j)\,\zeta(2k - 2j)
        \qquad (k \geq 2),
        \]
        and check it computes $\zeta(4)$ from $\zeta(2)$ and
        $\zeta(6)$ from $\zeta(2), \zeta(4)$ --- all even
        zeta values from the single seed $\frac{\pi^2}6$,
        with no new integration.
\end{enumerate}
\end{problem}
